\documentclass[12pt]{article}
\usepackage[T1]{fontenc}
\usepackage[utf8]{inputenc}
\usepackage{lmodern}
\usepackage{textcomp}
\usepackage{enumitem}
\usepackage{geometry}
\usepackage{graphicx}
\usepackage{amsmath, amssymb}
\geometry{margin=1in}
\begin{document}
\begin{center}
Biology - Undergraduate Level Assessment 23 \
Total Marks: 40 \
Answer all questions.
\end{center}

\section*{Multiple Choice (10 marks)}
\begin{enumerate}[label=\arabic*.]
\item What unusual feature of the cycads makes them are remarkable group among the seed plants?
\begin{enumerate}[label=(\alph*)]
    \item They have flowers that bloom for several years before producing seeds
    \item They have a unique leaf structure
    \item They can survive in extremely cold climates
    \item They produce large seeds
    \item Their seeds can remain dormant for up to a century before germinating
\end{enumerate}
\item Describe the function of the lateral-line system in fishes.
\begin{enumerate}[label=(\alph*)]
    \item The lateral-line system in fishes helps with the digestion of food.
    \item The lateral-line system in fishes is used for excreting waste materials.
    \item The lateral-line system in fishes is involved in the secretion of hormones.
    \item The lateral-line system in fishes is used for breathing.
    \item The lateral-line system in fishes is used for vision.
\end{enumerate}
\item Which of the following would maintain a constant body temperature in spite of changes in the environmental temperature: frog, robin, fish, dog, or lizard?
\begin{enumerate}[label=(\alph*)]
    \item dog, robin, fish
    \item robin, dog
    \item lizard, frog
    \item fish, robin
    \item fish, frog
\end{enumerate}
\item Compare the methods of obtaining oxygen in the earthworm, the frog, the fish and man.
\begin{enumerate}[label=(\alph*)]
    \item Earthworms use gills, frogs use both gills and lungs, fish use lungs, and humans use skin.
    \item Earthworms and humans use lungs, frogs and fish use gills.
    \item All use gills for obtaining oxygen.
    \item Earthworms and humans use gills, frogs use skin, and fish use lungs.
    \item Earthworms use lungs, frogs use gills, fish use skin, and humans use gills.
\end{enumerate}
\item Explain how an estimate of the age of a rock is made on the basisof the radioactive elements present.
\begin{enumerate}[label=(\alph*)]
    \item By assessing the fossil content found within the rock layers
    \item By comparing the density of the rock to a known age database
    \item By examining the color and texture of the rock
    \item By counting the layers of the rock
    \item By measuring the amount of carbon-14 remaining in the rock
\end{enumerate}
\end{enumerate}

\section*{True / False (10 marks)}
\textit{Answer True or False}
\begin{enumerate}[label=\arabic*.]
\setcounter{enumi}{5}
\item True or False: Insulin decreases blood sugar levels by promoting the breakdown of glycogen in the liver.
\item True or False: Sexual dimorphism is primarily caused by artificial selection, which is driven by human preference for specific traits.
\item True or False: Adenosine triphosphate functions primarily as an energy currency and does not act as a second messenger for calcium release.
\item True or False: Complete equilibrium in a gene pool is not expected due to ongoing evolutionary processes and environmental changes.
\item True or False: An increase in solar radiation is a biotic factor that enhances the growth rate of all populations in an ecosystem.
\end{enumerate}

\section*{Long Form Response (20 marks)}
\textit{Answer each question with a clear and complete explanation.}
\begin{enumerate}[label=\arabic*.]
\setcounter{enumi}{10}
\item Discuss how natural selection may lead to the rapid emergence of drug-resistant HIV strains following treatment with 3 TC, including the role of pre-existing viral populations.
\item Discuss the evolutionary advancements found in Selaginella that differentiate it from ferns, focusing on vascular structures and reproductive strategies.
\end{enumerate}

\vfill
\noindent\textit{End of Paper}
\end{document}